% 从项目根目录编译：latexmk -xelatex -outdir=build/quick quick-guide.tex
\documentclass[language=chinese,margin=right,11pt]{lumenbook}
\geometry{left=22mm,right=22mm,top=18mm,bottom=18mm,marginparwidth=0pt,marginparsep=0pt}
\title{生成模型统一框架速览}
\author{}
\date{}
\pagestyle{plain}
\setlength{\parindent}{0pt}
\setlength{\parskip}{5pt}
\begin{document}
\begin{center}
{\LARGE\bfseries 生成模型统一框架速览}\par
{\small 第 5 部分的快速阅读版，与全书第 0--5 章配套}
\end{center}

\section*{从一个概率演化问题出发}

所有模型先写清五个要素，再确定训练和生成算法。
\[
 \boxed{\text{端点分布}\ \longrightarrow\ \text{概率路径与耦合}
 \ \longrightarrow\ \text{可学习的场}\ \longrightarrow\ \text{生成动力学}}
\]
训练把场估计出来，生成把初始随机样本沿动力学推进。
误差分析连接这两个环节，并检查最终分布是否接近数据目标。

\par\noindent
\begin{tabularx}{\textwidth}{@{}llX@{}}
\toprule
要素 & 记号 & 必须说明的内容 \\
\midrule
数据与先验 & \(D\sim\mu_0,\ N\sim\mu_1\) & 数据尺度、先验类型、实际噪声端是否等于所选先验 \\
端点耦合 & \((D,N)\sim\pi\) & 独立、OT、教师生成或其他联合律 \\
概率路径 & \(p_t\) & 每个时刻的分布，可由扰动核、插值或 SDE 定义 \\
学习对象 & \(s_t,u_t,f_t\) & Score、速度、漂移或可换算的噪声／去噪预测 \\
生成算法 & \(\rho_\tau,\ v_\tau,b_\tau\) & 初始分布、时间方向、ODE／SDE 和数值网格 \\
\bottomrule
\end{tabularx}

\subsection*{先固定时间方向}
\[
 \underbrace{t=0\ \text{数据}\ \longrightarrow\ t=1\ \text{噪声}}_{\text{加噪时间}},
 \qquad \tau=1-t,\qquad
 \underbrace{\tau=0\ \text{噪声}\ \longrightarrow\ \tau=1\ \text{数据}}_{\text{生成时间}}.
\]
\(\rho_\tau=p_{1-\tau}\)，加噪速度 \(u_t\) 与生成速度满足
\(v_\tau=-u_{1-\tau}\)。空间 Score 为
\(s_t(x)=\nabla_x\log p_t(x)\)，在生成时间中按 \(s_{1-\tau}(x)\) 求值。

\subsection*{两条密度方程构成共同基础}
确定性输运使用连续性方程
\[
 \dot X_t=u_t(X_t),\qquad \partial_t p_t=-\nabla\cdot(p_tu_t).
\]
随机扩散使用 It\^o SDE 及 Fokker--Planck 方程
\[
 \dd X_t=f_t(X_t)\dd t+G_t(X_t)\dd W_t,\qquad a_t=G_tG_t^\top,
\]
\[
 \partial_t p_t=-\nabla\cdot(f_tp_t)+\tfrac12\sum_{ij}\partial_{ij}(a_{ij,t}p_t).
\]
在内部正则区间，经典密度推导要求光滑性、可积性和相应边界条件。
奇异端点可以通过测度弱形式或适当极限处理。

\clearpage
\section*{训练统一为条件回归，模型区别落在路径选择}

固定时刻，设 \(I_t\) 为可采样中间状态，\(T_t\) 为平方可积的可计算标签。平方回归满足
\[
 \mathbb E\norm{F_\theta(I_t,t)-T_t}^2
 =\mathbb E\norm{F_\theta(I_t,t)-\mathbb E[T_t\mid I_t]}^2+C_t.
\]
\(C_t\) 是标签的条件剩余方差，与网络参数无关。
因此总体最优输出是条件均值。生成误差需要控制超出该常数的场误差，
再结合模型容量、优化和统计泛化条件。

\par\noindent
\begin{tabularx}{\textwidth}{@{}llX@{}}
\toprule
训练方式 & 可采样标签 & 总体最优输出 \\
\midrule
噪声回归 & \(\varepsilon\) & \(\mathbb E[\varepsilon\mid X_t=x]\) \\
去噪回归 & \(D\) & \(\mathbb E[D\mid X_t=x]\) \\
DSM & \(\nabla_x\log q_t(X_t\mid D)\) & \(s_t(x)\)，条件核求导可与边缘积分交换 \\
CFM / 随机插值 & \(\dot I_t\) & \(u_t(x)=\mathbb E[\dot I_t\mid I_t=x]\) \\
桥迭代回归 & 指定方向的漂移监督 & 对应桥迭代所需的场，另需端点投影与迭代条件 \\
\bottomrule
\end{tabularx}

\subsection*{把八个主线对象放入同一张表}
\par\noindent
\begin{tabularx}{\textwidth}{@{}llX@{}}
\toprule
模型或框架 & 路径如何确定 & 学习与生成的对应关系 \\
\midrule
DDPM & 指定离散前向链 & 学习反向核参数，生成时逐步去噪 \\
DDIM & 保留指定单时刻边缘 & 复用去噪预测，选择隐式递推与随机强度 \\
Score SDE & 指定前向扩散 & 学习边缘 Score，生成解反向 SDE 或概率流 ODE \\
CNF & 参数化连续流 & 用变量变换计算似然，生成时解 ODE \\
FM / CFM & 指定条件概率路径 & 回归边缘速度，生成时解相应 ODE \\
Rectified Flow & 指定线性插值与耦合 & 回归端点位移；Reflow 重新选择训练耦合 \\
Schrödinger Bridge & 参考律与两端约束 & 优化路径 KL，学习控制漂移或双势 \\
随机插值 & 端点、插值及辅助噪声 & 回归速度／Score，构造共享边缘的 ODE／SDE \\
\bottomrule
\end{tabularx}

经典 CNF 的似然训练使用密度变化公式。
FM 提供通过条件回归学习连续流的另一条训练路线。
桥问题的共同回归步骤还需放回其路径空间优化目标中解释。
固定边缘路径时可转换场，改变耦合或路径时需要重新确定训练标签。

\clearpage
\section*{统一的场转换公式}

在 \(p_t>0\) 的正则内部时段，令 \(s_t=\nabla\log p_t\)。
一般扩散的概率流速度、反向生成漂移与生成 ODE 速度满足
\[
 \boxed{u_t=f_t-\tfrac12\operatorname{div}a_t-\tfrac12a_ts_t},
 \qquad (\operatorname{div}a_t)_i=\sum_j\partial_j a_{ij,t},
\]
\[
 \boxed{b_\tau=-f_t+\operatorname{div}a_t+a_ts_t},\qquad
 \boxed{v_\tau=-u_t},\quad t=1-\tau.
\]
若 \(a_t=g(t)^2I_d\)，得到 \(u_t=f_t-g^2s_t/2\)、
\(b_\tau=-f_t+g^2s_t\)、\(v_\tau=-f_t+g^2s_t/2\)。
状态依赖扩散保留矩阵散度项。

\subsection*{高斯路径把常见网络输出接起来}
设 \(X_t=\alpha_tD+\sigma_t\varepsilon\)，\(\varepsilon\sim\mathcal N(0,I_d)\)
与 \(D\) 独立，记 \(e_t^*=\mathbb E[\varepsilon\mid X_t=x]\)、
\(d_t^*=\mathbb E[D\mid X_t=x]\)。
\par\noindent
\begin{tabularx}{\textwidth}{@{}llX@{}}
\toprule
已知量 & 目标量 & 公式与使用条件 \\
\midrule
噪声均值 & Score & \(s_t=-e_t^*/\sigma_t\)，\(\sigma_t>0\) \\
Score & 干净均值 & \(d_t^*=(x+\sigma_t^2s_t)/\alpha_t\)，\(\alpha_t\ne0\) \\
干净均值与噪声均值 & 路径速度 & \(u_t=\dot\alpha_td_t^*+\dot\sigma_te_t^*\) \\
Score & 路径速度 & \(u_t=(\dot\alpha_t/\alpha_t)x+K_ts_t\)，\(\alpha_t\ne0\) \\
路径速度 & Score & \(s_t=(\dot\alpha_tx-\alpha_tu_t)/(\sigma_t\Delta_t)\)，\(\sigma_t\Delta_t\ne0\) \\
\bottomrule
\end{tabularx}
\[
 K_t=(\dot\alpha_t/\alpha_t)\sigma_t^2-\sigma_t\dot\sigma_t,
 \qquad \Delta_t=\alpha_t\dot\sigma_t-\sigma_t\dot\alpha_t.
\]
分母很小时误差会放大。相关端点、矩阵噪声或其他扰动核需按实际条件律换算。

\subsection*{线性路径的一个直接用法}
若 \(X_t=(1-t)D+t\varepsilon\)，则 \(\Delta_t=1\)，内部时刻有
\[
 u_t=-\frac{x+ts_t}{1-t},\qquad
 s_t=-\frac{x+(1-t)u_t}{t},\qquad v_\tau=-u_{1-\tau}.
\]
扩散 \(v\)-prediction 标签 \(V^{\mathrm{diff}}=\alpha\varepsilon-\sigma D\)
还需结合路径参数解释。在 \(\alpha=\cos\vartheta,\sigma=\sin\vartheta\) 时，
\(u_t=\vartheta'(t)\mathbb E[V_t^{\mathrm{diff}}\mid X_t=x]\)。

\clearpage
\section*{生成统一为初值问题，采样器负责离散求解}

先指定起点 \(Y_0\sim p_1\)，再在递增的生成时间 \(\tau\) 中求解。
有限噪声端用其他先验替代 \(p_1\) 时，应记录初始分布失配。

\par\noindent
\begin{tabularx}{\textwidth}{@{}llX@{}}
\toprule
生成方式 & 连续动力学 & 一个常用离散步 \\
\midrule
确定性流 & \(\dot Z=v_\tau(Z)\) & Euler，\(Z_{n+1}=Z_n+h_nv_\theta(Z_n,\tau_n)\) \\
随机扩散 & \(\dd Y=b_\tau\dd\tau+G\dd W\) & Euler--Maruyama，\(Y_{n+1}=Y_n+h_nb_\theta+\sqrt{h_n}G\xi_n\) \\
随机插值扩散 & \(b_\tau=v_\tau+\kappa s_{1-\tau}\) & 噪声尺度为 \(\sqrt{2\kappa h_n}\)，\(\kappa\) 仅依赖时间 \\
\bottomrule
\end{tabularx}
其中 \(\xi_n\) 是独立标准高斯，所有场在当前状态与正确时刻求值。
Heun 每步通常需两次场评估，RK4 通常需四次。
随机 Predictor--Corrector 还需计入固定噪声时刻的内部校正调用。

\subsection*{DDPM 与 DDIM 的离散版本放在哪里}
令 \(0<\beta_k<1\)、\(r_k=1-\beta_k\)、\(A_k^2=\prod_{j=1}^kr_j\)、\(S_k^2=1-A_k^2\)。
DDPM 用当前噪声预测确定反向均值
\[
 Y_{k-1}=\frac{Y_k-\beta_k\varepsilon_\theta(Y_k,k)/S_k}{\sqrt{r_k}}
                  +\sqrt{\nu_k}\xi_k.
\]
DDIM 先得到 \(\widehat D_k=(Y_k-S_k\varepsilon_\theta)/A_k\)，再跨到 \(j<k\)
\[
 Y_j=A_j\widehat D_k+\sqrt{S_j^2-\delta_{k,j}^2}\,\varepsilon_\theta
                           +\delta_{k,j}\xi,
 \qquad
 \delta_{k,j}^2=\eta^2\frac{S_j^2}{S_k^2}\left(1-\frac{A_k^2}{A_j^2}\right).
\]
常用 \(0\le\eta\le1\)，\(\eta=0\) 给出确定性更新。用累计系数处理跳步，有限步误差随网格改变。

\subsection*{相同边缘的条件}
精确场、匹配初始律和密度演化唯一性使概率流 ODE 与对应 SDE 共享边缘。
它们仍可具有不同的轨迹、端点耦合、二次变差和时间相关性。
训练时的条件插值线段与积分边缘速度产生的轨迹也需分别计算。
路径 KL 最优性、OT 最优性和边缘匹配各自对应所定义的优化条件。

\clearpage
\section*{用一张误差预算检查理论与实现}

将数值输出律记为 \(\widehat\rho_1^h\)，设计的平滑目标为 \(\rho_1^\epsilon\)。
\[
 \begin{aligned}
 W_2(\widehat\rho_1^h,\mu_0)
 &\le\underbrace{W_2(\widehat\rho_1^h,\widehat\rho_1)}_{\text{数值离散}}+
       \underbrace{W_2(\widehat\rho_1,\rho_1^\epsilon)}_{\text{场误差与起点失配}}\\
 &\quad+\underbrace{W_2(\rho_1^\epsilon,\mu_0)}_{\text{端点平滑／截断}}.
 \end{aligned}
\]
学习速度的空间 Lipschitz 常数为 \(L\)、参考边缘上场误差范数为 \(e(\tau)\) 时，
\[
 W_2(\widehat\rho_1,\rho_1)
 \le e^L W_2(\widehat\rho_0,\rho_0)+\int_0^1e^{L(1-\tau)}e(\tau)\dd\tau.
\]
该界要求相关解适定且二阶矩有限。回归风险需要经时间权重换算才能控制 \(e\)。

对同一非退化扩散的两个路径律，在 Girsanov 条件下
\[
 D_{\mathrm{KL}}(P\|\widehat P)
 =D_{\mathrm{KL}}(P_0\|\widehat P_0)
 +\tfrac12\mathbb E_P\int_0^1\norm{a_\tau^{-1/2}(b_\tau-\widehat b_\tau)}^2\dd\tau.
\]
终点 KL 由路径 KL 上界控制。选择 Wasserstein、KL 或全变差时，需检查各自的矩、
绝对连续性和端点条件。

\par\noindent
\begin{tabularx}{\textwidth}{@{}llX@{}}
\toprule
环节 & 主要条件 & 实际核对量 \\
\midrule
路径与标签 & 联合律正确、标签可积 & 条件均值、后验权重、端点与时间抽样 \\
ODE & 空间 Lipschitz 与增长控制等 & 初值稳定性、爆破、端点系数 \\
SDE & 相应存在唯一性与矩条件 & 漂移、协方差、状态依赖修正 \\
数值求解 & 求解器阶数与正则性相容 & 网格细化、NFE、均值／协方差、误差容限 \\
条件与物理结构 & 正确条件律、可行状态空间 & 条件起点、似然、约束残差和边界通量 \\
\bottomrule
\end{tabularx}

\subsection*{一个能贯穿全部步骤的核对模型}
取 \(D\sim\mathcal N(0,cI_d)\)、\(N\sim\mathcal N(0,I_d)\) 独立，
\(I_t=(1-t)D+tN\)，\(H_t=c(1-t)^2+t^2\)。则
\[
 s_t=-x/H_t,\qquad u_t=H_t'x/(2H_t),\qquad
 Z_\tau=\sqrt{H_{1-\tau}}\,Z_0,\quad Z_0\sim\mathcal N(0,I_d).
\]
这组公式可同时检查训练标签、场转换、时间方向和生成方差。

\clearpage
\section*{把统一框架用到条件生成与新模型}

\subsection*{条件生成仍沿用同一个框架}
把 \(p_t\) 换成 \(p_t(\cdot\mid c)\)，对相同条件联合律定义路径与回归标签。
生成起点应匹配 \(p_1(\cdot\mid c)\)。
对逆问题 \(y=\mathcal A(D)+\eta\)，精确带噪后验 Score 为
\[
 s_t(x\mid y)=s_t(x)+\nabla_x\log\int p(y\mid d)\,p_0(\dd d\mid X_t=x).
\]
用去噪均值代入干净似然需要分析所引入的近似。
CFG 在本书的尺度约定下取
\(s^{(w)}=s+w(s(\cdot\mid c)-s)\)，\(w=1\) 为普通条件 Score。
一般 guidance 放大后的终点律需结合整个采样动力学验证。

\subsection*{物理约束先落实到状态表示}
对线性约束 \(Az=c\)，写成 \(z=z_*(c)+Q\xi\)，其中 \(AQ=0,Q^\top Q=I\)。
在独立坐标 \(\xi\) 上设计先验、插值和动力学，可保持该约束。
非线性约束 \(C(z)=0\) 下，It\^o SDE 还需满足
\[
 \nabla C\,G=0,\qquad \nabla C\cdot f+\tfrac12a:\nabla^2C=0.
\]
流场可选无散基或 Fourier Leray 投影；涡度场还需核对零均值与速度恢复算子。
物理时间、生成时间、网格加密与不确定性覆盖分别记录。

\subsection*{阅读或提出一个模型时，依次回答六个问题}
\begin{enumerate}
 \item \textbf{目标是什么？} 写出数据、先验、条件信息及状态空间。
 \item \textbf{概率路径怎么来？} 写出扰动核、端点耦合、插值或参考扩散。
 \item \textbf{监督标签能否直接采样？} 推导条件均值及相应总体风险。
 \item \textbf{生成解什么方程？} 固定时间方向、输出换算和初始分布。
 \item \textbf{数学上是否适定？} 检查正则性、端点、边界、约束及必要的引用定理。
 \item \textbf{误差如何验证？} 分析场误差、起点、离散与端点误差，并做解析算例和网格核对。
\end{enumerate}

\subsection*{需要展开时回到哪一章}
\par\noindent
\begin{tabularx}{\textwidth}{@{}llX@{}}
\toprule
阅读目标 & 对应章节 & 使用方式 \\
\midrule
数学与记号 & 第 0 章 & 概率、输运、SDE、时间规范 \\
模型全流程 & 第 1 章 & 查训练与生成各自做什么 \\
训练目标 & 第 2 章 & 查 ELBO、DSM、CFM 与桥优化推导 \\
生成算法 & 第 3 章 & 查离散递推、ODE／SDE 与数值误差 \\
输出转换 & 第 4 章 & 查噪声、Score、速度、漂移及退化条件 \\
理论与应用 & 第 5 章 & 查统一、适定性、误差、约束与后验生成 \\
\bottomrule
\end{tabularx}

\subsection*{原始文献与阅读边界}
{\small
\href{https://arxiv.org/abs/2006.11239}{DDPM}；
\href{https://arxiv.org/abs/2010.02502}{DDIM}；
\href{https://arxiv.org/abs/2011.13456}{Score SDE}；
\href{https://arxiv.org/abs/1806.07366}{CNF / Neural ODE}；
\href{https://arxiv.org/abs/2210.02747}{Flow Matching}；
\href{https://arxiv.org/abs/2209.03003}{Rectified Flow}；
\href{https://arxiv.org/abs/1308.0215}{Schrödinger Bridge}；
\href{https://jmlr.org/papers/v26/23-1605.html}{Stochastic Interpolants}。
详细文献、推导和假设见全书。本速览为快速阅读版，独立数学审阅待完成。}
\end{document}
